\newpage
\phantomsection
\label{prf:darboux-refinement-squeeze}
\LRAProofFor{lem:darboux-refinement-squeeze}
\begin{remark*}[Return]\hyperref[lem:darboux-refinement-squeeze]{Return to Lemma}\end{remark*}
\begin{lemma*}[Darboux Sums Under Refinement]
Refinement raises lower sums and lowers upper sums.
\end{lemma*}
\begin{proof}[Professional Standard Proof]
\LRAProofBodyStart
It suffices to add one point \(c\in(x_{k-1},x_k)\). Let
\[
  I=[x_{k-1},x_k],\qquad I_1=[x_{k-1},c],\qquad I_2=[c,x_k].
\]
Write \(m,m_1,m_2\) for the infima of \(f\) on \(I,I_1,I_2\), and
\(M,M_1,M_2\) for the corresponding suprema. Since \(I_1,I_2\subset I\),
\[
  m\le m_1,\quad m\le m_2,\qquad M_1\le M,\quad M_2\le M.
\]
Therefore
\[
  m(c-x_{k-1})+m(x_k-c)
  \le m_1(c-x_{k-1})+m_2(x_k-c),
\]
so the lower sum rises. Similarly,
\[
  M_1(c-x_{k-1})+M_2(x_k-c)
  \le M(c-x_{k-1})+M(x_k-c),
\]
so the upper sum falls. Repeating for each added point proves the chain.
\end{proof}
\begin{proof}[Detailed Learning Proof]
\LRAProofBodyStart
When a slab is split, the new lower rectangles may use larger infima because
each smaller interval has fewer points where the function can dip. The new
upper rectangles may use smaller suprema for the same reason. This is the
monotone control Darboux extracts from refinement.
\end{proof}
\begin{remark*}[Proof structure]
Check the effect of adding one partition point, then iterate.
\end{remark*}
\begin{dependencies}
\begin{itemize}
  \item \hyperref[def:darboux-sums]{Lower and Upper Darboux Sums}.
  \item \hyperref[def:partition-refinement]{Refinement of a Partition}.
\end{itemize}
\end{dependencies}
\clearpage
